% Part 2 – Theory of Energy–Based Models

% ---------------------------------------------------------------------------
% This section delves into the theoretical foundations of energy–based models
% (EBMs).  It expands upon the brief summary provided in Part 1 by defining
% energy functions, partition functions, and score functions, deriving the
% log–likelihood gradient used for training, and analysing sampling schemes
% such as Langevin dynamics and rejection sampling.  We also revisit the
% mixture–of–Gaussians example to illustrate how the score function can
% ignore mixing weights, leading to multimodal sampling difficulties.

\section{Energy--Based Models: Theory}

\subsection{Definition of an Energy Function}

An \emph{energy–based model} specifies a distribution over a high–
dimensional variable $\mathbf{x}$ through a scalar–valued \emph{energy
function} $E_\theta(\mathbf{x})$, parameterised by $\theta$:
\begin{equation}
    p_\theta(\mathbf{x}) = \frac{\exp\big(-E_\theta(\mathbf{x})\big)}{Z_\theta},
    \quad \text{with} \quad Z_\theta = \int \exp\big(-E_\theta(\mathbf{x})\big)
    \,\mathrm{d}\mathbf{x}.
    \label{eq:ebm-density}
\end{equation}
Here $Z_\theta$ is the \emph{partition function} that normalises the density.
In contrast to probabilistic models that explicitly compute $Z_\theta$, EBMs
are often trained by maximising the log–likelihood of the data without
evaluating $Z_\theta$ exactly.  As emphasised by LeCun et~al., the absence
of a normalisation requirement confers great flexibility: probabilistic
models must be properly normalised, which can require evaluating
intractable integrals over all possible configurations, whereas EBMs have
no such requirement and thus circumvent the problem of normalisation
entirely【861141214805006†L26-L33】.  This property allows the design of
loss functions that encourage the energy assigned to observed data to be
smaller than that assigned to unobserved configurations【861141214805006†L46-L56】.

Equation~\eqref{eq:ebm-density} suggests that low–energy points have high
probability under the model.  Learning and inference are therefore
formulated in terms of minimising $E_\theta(\mathbf{x})$ with respect to
\emph{unobserved} variables and adjusting $\theta$ so that observed
configurations lie in low–energy regions【861141214805006†L46-L63】.  We
distinguish between two sets of variables: observed inputs $\mathbf{x}$ and
variables to be predicted $\mathbf{y}$; the energy function $E(\mathbf{y},
\mathbf{x})$ measures their compatibility, and inference corresponds to
finding the value of $\mathbf{y}$ that minimises $E(\mathbf{y},\mathbf{x})$ at
fixed $\mathbf{x}$【861141214805006†L261-L267】.

\subsection{Partition Function and Log–Likelihood Gradient}

Although the density in Eq.~\eqref{eq:ebm-density} depends on $Z_\theta$, we
often avoid computing the partition function explicitly by working with
gradients of the log–likelihood.  Given a dataset $\{\mathbf{x}^{(n)}\}$, the
negative log–likelihood loss is
\begin{equation}
    \mathcal{L}(\theta) = -\frac{1}{N}\sum_{n=1}^N \log p_\theta\big(\mathbf{x}^{(n)}\big)
    = \frac{1}{N}\sum_{n=1}^N \Big(E_\theta\big(\mathbf{x}^{(n)}\big) + \log Z_\theta\Big).
\end{equation}
The gradient of $\mathcal{L}(\theta)$ with respect to $\theta$ can be derived
using the log–likelihood trick:
\begin{align}
    \nabla_\theta \log p_\theta(\mathbf{x})
    &= \nabla_\theta \big(-E_\theta(\mathbf{x}) - \log Z_\theta\big) \\
    &= -\nabla_\theta E_\theta(\mathbf{x}) + \frac{1}{Z_\theta} \nabla_\theta Z_\theta
      \quad (\text{by differentiating under the integral sign}) \\
    &= -\nabla_\theta E_\theta(\mathbf{x}) + \int \exp\big(-E_\theta(\mathbf{x}')\big)
      \nabla_\theta E_\theta(\mathbf{x}') \; p_\theta(\mathbf{x}') \,\mathrm{d}\mathbf{x}' \\
    &= -\nabla_\theta E_\theta(\mathbf{x}) + \mathbb{E}_{\mathbf{x}' \sim p_\theta}
      \big[\nabla_\theta E_\theta(\mathbf{x}')\big],
    \label{eq:ebm-gradient}
\end{align}
where we have used the fact that $\nabla_\theta Z_\theta = \int 
\exp\big(-E_\theta(\mathbf{x}')\big)(-\nabla_\theta E_\theta(\mathbf{x}'))\,\mathrm{d}\mathbf{x}'$
and that $p_\theta(\mathbf{x}') = \exp\big(-E_\theta(\mathbf{x}')\big)/Z_\theta$.
Equation~\eqref{eq:ebm-gradient} is the key result underpinning maximum
likelihood training for EBMs: the gradient consists of a \emph{positive
phase} term $-\nabla_\theta E_\theta(\mathbf{x})$ that pushes down the energy
of real data and a \emph{negative phase} term $\mathbb{E}_{p_\theta}[\nabla_\theta
E_\theta(\mathbf{x}')]$ that pushes up the energy of model samples.  The
training objective therefore encourages lower energy for data points and
higher energy for synthetic samples drawn from the current model.

\subsection{Sampling and the Challenge of the Negative Phase}

The main computational difficulty in evaluating Eq.~\eqref{eq:ebm-gradient}
lies in estimating the negative phase expectation over $p_\theta$.  Since
directly sampling from $p_\theta$ is generally intractable, one must resort
to approximate methods.  \emph{Contrastive divergence} (CD) is a popular
technique introduced by Hinton that replaces the expectation over the model
distribution with a short Markov chain starting from the data.  In CD-$k$
training, for each data point $\mathbf{x}$, one initialises a chain at
$\mathbf{x}$, performs $k$ steps of a Markov transition kernel (often
Langevin dynamics or Gibbs sampling), and uses the resulting sample
$\tilde{\mathbf{x}}$ to approximate the negative phase gradient.  The
update direction is then
\begin{equation}
    \nabla_\theta \mathcal{L}_{\text{CD}} \approx -\frac{1}{N}\sum_{n=1}^N
    \nabla_\theta E_\theta\big(\mathbf{x}^{(n)}\big) +
    \frac{1}{N}\sum_{n=1}^N \nabla_\theta E_\theta\big(\tilde{\mathbf{x}}^{(n)}\big),
\end{equation}
where $\tilde{\mathbf{x}}^{(n)}$ denotes the sample after $k$ Markov steps.  The
approximation becomes exact as $k\to\infty$ and the chain equilibrates, but
in practice very few steps (e.g., $k=1$ or $k=10$) often suffice to
produce useful updates.

\subsection{Langevin Dynamics for Sampling}

One of the most widely used sampling techniques for EBMs is \emph{Langevin
dynamics}.  It arises from the discretisation of a stochastic differential
equation (SDE) that simulates Brownian motion on an energy landscape.
Starting from an initial configuration $\mathbf{x}_0$, Langevin updates are
given by
\begin{equation}
    \mathbf{x}_{t+1} = \mathbf{x}_t - \frac{\eta}{2}
    \nabla_{\mathbf{x}} E_\theta(\mathbf{x}_t) + \sqrt{\eta}\,\boldsymbol{\xi}_t,
    \quad \boldsymbol{\xi}_t \sim \mathcal{N}(\mathbf{0}, \mathbf{I}),
    \label{eq:langevin}
\end{equation}
where $\eta$ is the step size and $\boldsymbol{\xi}_t$ is a Gaussian noise
vector.  Equation~\eqref{eq:langevin} can be derived by discretising the
overdamped Langevin SDE $\mathrm{d}\mathbf{x} = -\nabla_{\mathbf{x}}
E_\theta(\mathbf{x})\,\mathrm{d}t + \sqrt{2}\,\mathrm{d}\mathbf{W}(t)$, where
$\mathbf{W}(t)$ is a Wiener process.  Under suitable conditions on
$E_\theta$, the Markov chain defined by Eq.~\eqref{eq:langevin} has
$p_\theta$ as its stationary distribution.  In practice, one clamps the
iterate to a valid range (e.g., $[0,1]$ for image pixels) after each step
to ensure that generated samples remain within the domain.

Langevin dynamics is used both during training (to generate negative phase
samples) and at inference time (to produce novel images or denoise noisy
images).  The quality of the samples depends on the number of steps, the
step size, and the balance between exploration (due to noise) and
exploitation (due to the gradient term).  Longer chains allow the Markov
process to mix more thoroughly but increase computational cost.

\subsection{Rejection Sampling and High--Dimensional Acceptance Rates}

As an alternative to gradient–based sampling, one could in principle use
\emph{rejection sampling} to draw from $p_\theta(\mathbf{x})$.  Suppose
one has a proposal distribution $q(\mathbf{x})$ from which samples are
readily drawn and a target density proportional to $\tilde{p}(\mathbf{x})
\propto \exp\big(-E_\theta(\mathbf{x})\big)$.  To perform rejection sampling
we must find a constant $M$ such that
\begin{equation}
    \tilde{p}(\mathbf{x}) \leq M\,q(\mathbf{x}) \quad \text{for all }\mathbf{x}.
\end{equation}
Given a sample $\mathbf{x}\sim q$ and a uniform random number $u\sim
\mathrm{Uniform}(0,M q(\mathbf{x}))$, the sample is accepted if $u \leq
\tilde{p}(\mathbf{x})$ and rejected otherwise.  The acceptance probability is
$1/M$, and the expected number of proposals before acceptance is $M$.

For high–dimensional data, however, the acceptance rate can be vanishingly
small.  Consider the case where both $q$ and $\tilde{p}$ are isotropic
Gaussian densities on $\mathbb{R}^d$ with zero mean but different standard
deviations $\sigma_q$ and $\sigma_p$.  The ratio of the unnormalised
density to the proposal density satisfies
\begin{equation}
    \frac{\tilde{p}(\mathbf{x})}{q(\mathbf{x})}
    = \Big(\frac{\sigma_q}{\sigma_p}\Big)^d \exp\Big(-\frac{\|\mathbf{x}\|^2}{2}
      \big(\sigma_p^{-2} - \sigma_q^{-2}\big)\Big) \leq
    \Big(\frac{\sigma_q}{\sigma_p}\Big)^d,
\end{equation}
so we may take $M=(\sigma_q/\sigma_p)^d$.  For MNIST, the dimensionality
is $d=784$.  If we choose $\sigma_q = 1.01\,\sigma_p$, then $M = 1.01^{784}
\approx \exp(784\times 0.00995) \approx \exp(7.8) \approx 2.44\times 10^3$.
Thus the acceptance rate is roughly $1/M \approx 4\times 10^{-4}$—one in
2,500 samples would be accepted on average.  This simple calculation
illustrates that rejection sampling is impractical for high–dimensional
problems like images.  In practice, gradient–based methods such as
Langevin dynamics or Hamiltonian Monte Carlo are preferred.

\subsection{Mixture of Gaussians and Score Independence from Mixing Weights}

Another instructive example concerns a mixture of two Gaussian densities
with disjoint supports.  Let $p_a(\mathbf{x})$ and $p_b(\mathbf{x})$ be
Gaussian densities with means $\boldsymbol{\mu}_a$, $\boldsymbol{\mu}_b$ and
covariances $\boldsymbol{\Sigma}_a$, $\boldsymbol{\Sigma}_b$, respectively,
such that $\text{supp}(p_a) \cap \text{supp}(p_b)=\varnothing$.  Consider the
mixture
\begin{equation}
    p(\mathbf{x}) = w_a p_a(\mathbf{x}) + w_b p_b(\mathbf{x}),
    \quad w_a, w_b \geq 0,\; w_a + w_b = 1.
\end{equation}
The log–density within the support of component $a$ is
\begin{equation}
    \log p(\mathbf{x}) = \log\big(w_a p_a(\mathbf{x})\big)
    = \log w_a + \log p_a(\mathbf{x}),
\end{equation}
and therefore
\begin{equation}
    \nabla_{\mathbf{x}} \log p(\mathbf{x}) = \nabla_{\mathbf{x}} \log p_a(\mathbf{x}),
    \quad \mathbf{x}\in\text{supp}(p_a).
\end{equation}
An analogous expression holds in $\text{supp}(p_b)$.  The gradient of the
log–density—\emph{the score function}—is thus independent of the mixing
weights $w_a$ and $w_b$ on either component.  The score is determined
solely by the local density and does not convey any information about the
global mixture proportions.  Consequently, when one performs Langevin
dynamics using the score, the chain will explore one component without
ever moving to the other.  The sampler cannot traverse from region $A$ to
region $B$ because the gradient vanishes on the boundary of the disjoint
supports.  This insight highlights the limitations of using simple gradient–
based dynamics for multimodal distributions.  In practise, one might
mitigate the issue by adding noise, tempering the energy landscape, or
employing non–local moves via importance resampling.

\subsection{Score Function and Independence from Normalisation}

It is often convenient to define the \emph{score function} of a model
$p_\theta(\mathbf{x})$ as the gradient of the log–density with respect to
$\mathbf{x}$:
\begin{equation}
    \mathbf{s}_\theta(\mathbf{x}) = \nabla_{\mathbf{x}} \log p_\theta(\mathbf{x})
    = -\nabla_{\mathbf{x}} E_\theta(\mathbf{x}) - \nabla_{\mathbf{x}} \log
    Z_\theta.
\end{equation}
Since $\log Z_\theta$ is constant with respect to $\mathbf{x}$, its gradient
vanishes.  Therefore
\begin{equation}
    \mathbf{s}_\theta(\mathbf{x}) = -\nabla_{\mathbf{x}} E_\theta(\mathbf{x}),
    \label{eq:score-from-energy}
\end{equation}
which shows that the score is independent of the partition function
$Z_\theta$.  This independence is a core advantage of working with scores:
one can learn gradients of log–densities without computing the normalising
constant.  In later sections we will see how this idea underpins
score–based generative models.

\subsection{Problems with Naïve Score Matching on Raw Data}

One might be tempted to directly match the model score function
$\mathbf{s}_\theta(\mathbf{x})$ to the true data score $\nabla_{\mathbf{x}} \log
p_{\text{data}}(\mathbf{x})$ by minimising the \emph{Fisher divergence},
\begin{equation}
    \mathcal{L}_{\text{SM}}(\theta) = \frac{1}{2}
    \int p_{\text{data}}(\mathbf{x}) \big\| \mathbf{s}_\theta(\mathbf{x}) -
    \nabla_{\mathbf{x}} \log p_{\text{data}}(\mathbf{x}) \big\|^2 \,
    \mathrm{d}\mathbf{x}.
\end{equation}
However, this loss involves computing the divergence of the score,
$\nabla\cdot\mathbf{s}_\theta(\mathbf{x})$, which is intractable in high
dimensions because it requires evaluating second derivatives of the model
output with respect to the input.  Moreover, the true data score
$\nabla_{\mathbf{x}} \log p_{\text{data}}(\mathbf{x})$ is unknown except in
simple parametric cases, making direct score matching impractical for
realistic datasets such as images.  These issues motivate alternative
training objectives, such as denoising score matching, which we discuss in
Part~5.

\subsection{Summary of Theoretical Insights}

In this part we introduced the core theoretical concepts underpinning
energy–based models.  We defined the energy function and its associated
unnormalised density, derived the maximum–likelihood gradient and
decomposed it into positive and negative phases, examined sampling via
Langevin dynamics and rejection sampling, and analysed the failure of the
score to incorporate mixture weights in multimodal settings.  We also
highlighted the independence of the score from the partition function and
discussed why naïvely matching scores on raw data is problematic in
practice.  These insights motivate the training strategies and design
decisions implemented in the next part of the report.